% Part: methods
% Chapter: proofs
% Section: reading-proofs

\documentclass[../../../include/open-logic-section]{subfiles}

\begin{document}

\olfileid{mth}{prf}{rea}
\olsection{నిరూపణలను చదవడం}

పాఠ్యపుస్తకాలు, వ్యాసాలలో కనిపించే నిరూపణలు మన ఉదాహరణల్లో ఇప్పటివరకు చూపిన అన్ని వివరాలను చాలా అరుదుగానే ఇస్తాయి. రచయితలు సందర్భాలను వేరుచూసినప్పుడు, పరోక్ష నిరూపణ ఇచ్చినప్పుడు తరచుగా ప్రత్యేకంగా చెప్పరు; నిర్వచనాన్ని ఉపయోగించిన సంగతి కూడా చెప్పకపోవచ్చు. కాబట్టి పాఠ్యపుస్తకంలో నిరూపణ చదివేటప్పుడు దాన్ని అర్థం చేసుకోవడానికి ఆ వివరాలను మీరే పూరించాల్సి రావచ్చు. నిరూపణలో చేయాల్సిన వివిధ కదలికలకు అలవాటు పడటానికీ ఇది మంచి సాధన. ఒక ఉదాహరణ చూద్దాం.

\begin{prop}[శోషణ]
అన్ని సమితులు $A$, $B$లకు,
\[
A \cap (A \cup B) = A
\]
\end{prop}

\begin{proof}
$z \in A \cap (A \cup B)$ అయితే $z \in A$; కాబట్టి $A \cap (A \cup B)
\subseteq A$. ఇప్పుడు $z \in A$ అనుకుందాం. అప్పుడు $z \in A \cup B$ కూడా; అందుచేత $z \in A \cap (A \cup B)$ కూడా.
\end{proof}

ముందరి శోషణ సూత్ర నిరూపణ చాలా సంక్షిప్తం. వాడిన నిర్వచనాలు ప్రస్తావించలేదు; నిరూపించాల్సింది ఏమిటో ముందుగా చెప్పలేదు; మొదలైనవి. దీన్ని విప్పి చూద్దాం. నిరూపించిన ప్రతిపాదన ఏ సమితులు $A$, $B$ల గురించిన సాధారణ వాదన; నిరూపణలో $A$ లేదా $B$ అన్నప్పుడు అవి ఏవైనా సమితులకు చరరాశులు. నిరూపణ స్థాపించే సాధారణ వాదనలు సమితుల తాదాత్మ్యానికి కావలసినవే: సమానత్వం ఎడమ వైపున ఉన్న ప్రతి \tetoken{మూలకం}{!!{element}}, కుడి వైపున \tetoken{ఒక మూలకం}{!!a{element}}; అలాగే వ్యతిరేక దిశలోనూ.

\begin{quote}
``$z \in A \cap (A \cup B)$ అయితే $z \in A$; కాబట్టి $A \cap (A \cup B)
  \subseteq A$.''
\end{quote}

ఇది సమానత్వ నిరూపణ మొదటి సగం: ఏదైనా~$z$ ఎడమ వైపున \tetoken{ఒక మూలకం}{!!a{element}} అయితే కుడి వైపున కూడా \tetoken{ఒక మూలకం}{!!a{element}} అని, అంటే $A \cap (A \cup B) \subseteq A$ అని, చూపుతుంది. $z \in A \cap (A \cup B)$ అనుకుందాం. రెండు సమితుల ఛేదనలో $z$ \tetoken{ఒక మూలకం}{!!{element}} అవడం అప్పుడూ అప్పుడే రెండు సమితులలోనూ \tetoken{ఒక మూలకం}{!!{element}} అవడం. కాబట్టి $z \in A$, అలాగే $z \in A \cup B$ అని నిగమించవచ్చు. ముఖ్యంగా $z \in A$; ఇదే చూపాల్సింది. మొదటి సగానికి అంతే చాలు కనుక నిరూపణలో మిగిలినది రెండవ సగం, అంటే $A \subseteq A \cap
(A \cup B)$ నిరూపణ అయి ఉండాలి.

\begin{quote}
``ఇప్పుడు $z \in A$ అనుకుందాం. అప్పుడు $z \in A \cup B$ కూడా; అందుచేత $z \in A \cap (A \cup B)$ కూడా.''
\end{quote}

$z \in A$ అని తీసుకొని మొదలుపెడతాం. ఎందుకంటే ఏ~$z$కైనా $z \in A$ అయితే $z \in A \cap (A \cup B)$ అని చూపుతున్నాం. $z
\in A \cap (A \cup B)$ అని చూపాలంటే ``$\cap$'' నిర్వచనం ప్రకారం (i) $z \in A$, (ii) $z \in A \cup B$ అని చూపాలి. ఇక్కడ (i) మన పరికల్పనే; దాన్ని మళ్లీ నిరూపించనవసరం లేదు, అందుకే నిరూపణ దాన్ని మళ్లీ ప్రస్తావించదు. (ii) కోసం, $z$ సమితుల సమ్మేళనంలో \tetoken{ఒక మూలకం}{!!a{element}} అప్పుడూ అప్పుడే ఆ సమితులలో కనీసం ఒకదానిలో \tetoken{మూలకం}{!!{element}} అని గుర్తుచేసుకోండి. $z \in A$ కాబట్టి, $A \cup B$ అనేది $A$, $B$ల సమ్మేళనం కాబట్టి ఇక్కడ అలా ఉంది. కనుక $z \in A \cup B$. (i) $z \in A$, (ii) $z \in A \cup B$ రెండూ చూపాం; అందువల్ల ``$\cap$'' నిర్వచనం ప్రకారం $z \in A \cap (A \cup B)$. నిరూపణ ఈ నిర్వచనాలను చెప్పదు; పాఠకుడికి అవి ఇప్పటికే తెలుసని భావిస్తుంది. తెలియకపోతే మళ్లీ గుర్తుచేసుకోవాలి. అప్పుడు $z
\in A$ నుంచి $z \in A \cup B$ ఎందుకు వస్తుందో, $z \in A$, $z \in A \cup B$ నుంచి $z \in A \cap (A \cup B)$ ఎందుకు వస్తుందో కూడా గుర్తించాలి.

పైనున్న నిరూపణలో ప్రతి విషయాన్ని స్పష్టంగా చెప్పిన మరో రూపం ఇది:
\begin{proof}{}
[సమితుల కోసం $=$ నిర్వచనం ప్రకారం $A \cap (A \cup B) = A$ అని చూపాలంటే (a) $A \cap (A \cup B) \subseteq A$, (b) $A \subseteq A \cap (A \cup B)$ అని చూపాలి. (a): $\subseteq$ నిర్వచనం ప్రకారం $z \in A \cap (A \cup B)$ అయితే $z \in A$ అని చూపాలి.] $z \in A
\cap (A \cup B)$ అయితే $z \in A$ [$\cap$ నిర్వచనం ప్రకారం $z
  \in A \cap (A \cup B)$ అప్పుడూ అప్పుడే $z \in A$ మరియు $z \in A \cup B$], కాబట్టి $A
\cap (A \cup B) \subseteq A$. [(b): $\subseteq$ నిర్వచనం ప్రకారం $z \in A$ అయితే $z \in A \cap (A \cup B)$ అని చూపాలి.] ఇప్పుడు [(1)] $z \in A$ అనుకుందాం. అప్పుడు [(2)] $z \in A \cup B$ [$z \in A$ లేదా $z \in B$ అని (1) నుంచి వస్తుంది; $\cup$ నిర్వచనం ప్రకారం అది $z
  \in A \cup B$], అందుచేత $z \in A \cap (A \cup B)$ కూడా [$\cap$ నిర్వచనం ప్రకారం $z \in A$, అంటే (1), మరియు $z
  \in A \cup B$, అంటే (2), కావాలి].
\sourcecorrection{OLTEMTHPRFREA-001}{మూలంలోని స్పష్టీకృత నిరూపణలో (b) కూడా పొరపాటున మొదటి ఉపసమితి దిశనే పునరావృతం చేసింది. సమానత్వానికి అవసరమైన వెనుక దిశ $A\subseteq A\cap(A\cup B)$గా సరిచేశాం.}
\sourcecorrection{OLTEMTHPRFREA-002}{మూల చివరి వివరణలో $z\in A\cup B)$ అని అదనపు మూసివేత కుండలీకరణం ఉంది. దాన్ని తొలగించాం.}
\end{proof}

\begin{prob}
$A \cup (A \cap B) = A$కు ఈ క్రింది నిరూపణను విస్తరించండి. అందులో వాడిన అన్ని నిగమన నమూనాలు, ప్రతి దశ పూర్వ పరికల్పనలు లేదా స్థాపించిన వాదనల నుంచి ఎందుకు వస్తుందో, ఏ నిర్వచనాలను ఎక్కడ ఉపయోగించాలో చెప్పండి.
\begin{proof}
  $z \in A \cup (A \cap B)$ అయితే $z \in A$ లేదా $z \in A \cap B$. $z \in A \cap B$ అయితే $z \in A$. ఏ $z \in A$ అయినా $\in A \cup (A
  \cap B)$ కూడా.
\end{proof}
\end{prob}

\end{document}
